% SWITCH-DAG-NROBP4
\section*{SWITCH-DAG-NROBP4: contacts as shared DAGs}

\textbf{STATUS: CANDIDATE with a complete proof in the declared passive monotone model; independent red-team review remains pending. P versus NP remains OPEN.}

\paragraph{Literature model.}
The requested contact-DAG model is a classical object.  An acyclic
switching-and-rectifier network is a directed acyclic graph with input-controlled switching
edges and free rectifier edges; an input is accepted when it enables an $s$--$t$ path.
The read-once version is commonly identified with a nondeterministic read-once branching
program (NROBP).  Rectifier networks themselves are DAGs whose source--sink reachability
relation is the computed object.

\paragraph{Declared model.}
Controlled contacts carry only positive clause-presence labels $x_c$.  Labels may repeat
arbitrarily, free routing and arbitrary DAG sharing are allowed, and correctness is required
only on the exact-weight-eight promise $D(n,8)$.  The charged resource $C$ is the number of
controlled contact occurrences.  Complement contacts, active relays, and retained sequential
state are outside this theorem.

Let
\[
M=8\binom n3
\]
and let $H_n$ be the clause-conflict graph.  Its independence number is exactly eight:
pairwise nonconflicting clauses have pairwise disjoint falsifying subcubes of size
$2^{n-3}$, so at most eight fit in the assignment cube, while all eight sign patterns on
one variable triple attain eight.

\paragraph{Cut-hitting lemma.}
Let $K$ be any controlled-contact $s$--$t$ cut and let $L(K)$ be its distinct clause labels.
Put $R=V(H_n)\setminus L(K)$.  Every promised input chosen only from $R$ disables the
entire cut.  If $|R|>8$, then $R$ contains a conflict pair because
$\alpha(H_n)=8$.  Adding any six further labels from $R$ produces a legal weight-eight
SAT input avoiding the cut, a contradiction.  Hence every controlled-contact cut contains
at least
\[
M-8
\]
distinct labels and therefore at least $M-8$ controlled edges.

Give controlled contacts unit capacity and free rectifier edges infinite capacity.  Directed
max-flow/min-cut yields at least $M-8$ $s$--$t$ paths disjoint in controlled contacts.
Every such path contains at least two controlled contacts: a one-contact path labelled
$x_c$ would accept an UNSAT exact cover containing $c$, for example the eight sign patterns
on the underlying variable triple.  Therefore
\[
C\ge2(M-8)=2M-16.
\]

\paragraph{Matching acyclic construction.}
Choose two distinct variable triples.  Let $A$ and $B$ be their respective eight-sign-pattern
exact covers.  Give $A$ source contacts only, $B$ sink contacts only, and every other clause
one source and one sink occurrence.  Orient each conflict pair from a source-capable endpoint
to a sink-capable endpoint and realize the oriented relation by free rectifier routing.
This is an acyclic switching-and-rectifier DAG using
\[
|A|+|B|+2(M-16)=2M-16
\]
controlled contacts.

Subject to independent review, this closes the passive monotone DAG contact complexity:
\[
\boxed{
C_{\mathrm{passive\ DAG}}(n,8)
=
2M-16
=
16\binom n3-16.
}
\]

In particular,
\[
C(4,8)=48,\qquad C(5,8)=144.
\]
This supersedes the weaker count-only $2\tau$ lower bound of SWITCH-CUT-MENGER3 for the
fixed-eight family.

\paragraph{D4 routing breadth.}
For the explicit 48-contact $D(4,8)$ DAG, the oriented source--sink conflict relation is a
$24\times24$ Boolean matrix with 96 ones.  Exhaustive maximal-biclique enumeration gives
72 maximal bicliques.  Both the fractional rectangle-cover LP and the exact integer cover
have value 20; the LP dual is integral and supplies a 20-edge fooling set.  Thus every
rectifier realization of this particular oriented relation has rectangle-cover cut breadth
at least 20.  This is a finite orientation-specific routing certificate, not a lower bound
for every possible optimal orientation.

\paragraph{What changed.}
Allowing contacts to act as DAGs does not invalidate the fixed-$m=8$ contact lower bound.
On the contrary, the structural cut argument survives arbitrary DAG sharing and arbitrary
positive-label repetition.  What remains uncharged is the free routing layer, whose direct
representation contains the conflict relation and can dominate total size.

\paragraph{Barrier.}
The model now coincides with a standard restricted complexity frontier: switching-and-rectifier
networks / nondeterministic branching programs.  Polynomial-size NBPs capture an
NL/poly-type nonuniform model, not all of polynomial time.  Thus even an asymptotically
strong NBP lower bound would still require a separate transfer theorem to constrain arbitrary
polynomial-time SAT algorithms.

\paragraph{Next hostile test.}
Allow normally-closed/complement contacts and hence full literal-labelled switch DAGs.
The cut-hitting proof is monotone and does not automatically survive negative contacts.
This is the most direct next test of whether exact-weight promise information collapses the
$2M-16$ contact requirement.

\begin{thebibliography}{9}
\bibitem{Jukna95}
S.~Jukna,
\emph{A Note on Read-$k$ Times Branching Programs},
RAIRO Theoretical Informatics and Applications 29(1):75--83, 1995.

\bibitem{Razborov91}
A.~A.~Razborov,
\emph{Lower Bounds for Deterministic and Nondeterministic Branching Programs},
FCT 1991, LNCS 529, 47--60.

\bibitem{IvanEtAl14}
S.~Iv\'an, A.~D.~Lelkes, J.~Nagy-Gy\"orgy, B.~Sz\"or\'enyi, and G.~Tur\'an,
\emph{Biclique Coverings, Rectifier Networks and the Cost of $\varepsilon$-Removal},
DCFS 2014, LNCS 8614, 174--185.

\bibitem{ChistikovEtAl24}
D.~Chistikov et al.,
discussion of NROBPs as acyclic read-once switching-and-rectifier networks,
ICDT 2024.
\end{thebibliography}

% END_OF_NOTE::SWITCH-DAG-NROBP4
